\documentclass{beamer}

\usepackage{framed,multicol}
\usepackage{color}
\usepackage{graphicx,caption,subcaption}
\usepackage{hyperref}
\hypersetup{
    colorlinks=true,
    linkcolor=blue,
    filecolor=green,
    urlcolor=cyan,
    citecolor=cyan,
}

\usetheme{default}
\usecolortheme{seahorse}
\usefonttheme{serif}
\setbeamertemplate{navigation symbols}{}
\setbeamertemplate{items}[ball]
\setbeamertemplate{section in toc}[ball unnumbered]
\setbeamertemplate{footline}[frame number]
\setbeamerfont{section in toc}{series=\bfseries}
\setbeamertemplate{part page}{
        \begin{beamercolorbox}[sep=8pt,center,wd=\textwidth]{part title}
            \bfseries\Large\insertpart\par
        \end{beamercolorbox}
}

\usepackage{unicode-math}
\unimathsetup{bold-style=ISO}
\setmathfont{XITS Math}
\usepackage{anyfontsize}
\usepackage{fontspec}
\setmainfont{Times New Roman}

\AtBeginDocument{
    \let\uglyepsilon\epsilon
    \renewcommand\epsilon{\varepsilon}
}
\usepackage{amsmath,envmath}
\usepackage{bm}
\renewcommand{\bm}{\symbf}

\usepackage[fontset=none]{ctex}
\setCJKmainfont{SimSun}[BoldFont=Noto Sans CJK SC Medium,RawFeature=+fwid,ItalicFont=KaiTi]
\setCJKsansfont{SimSun}[BoldFont=Noto Sans CJK SC Medium,RawFeature=+fwid,ItalicFont=KaiTi]
\setCJKmonofont{Noto Sans Mono CJK SC}[RawFeature=+fwid]
\newCJKfontfamily{\kaishu}{KaiTi}
\setbeamerfont{author}{family=\kaishu}
\setbeamerfont{date}{family=\kaishu}

\let\titleonly\title
\renewcommand{\title}[1]{\titleonly{\bfseries{#1}}}
\newcommand{\bib}[1]{\bibitem{bib:#1}}

\newenvironment{slide}{
    \begin{frame}{\bfseries{\sectiontitle}}
}{
    \end{frame}
}

\newenvironment{itemlist}{
    \begin{enumerate}
    \let\itemonly\item
    \renewcommand{\item}[1]{\itemonly \textbf{##1} \\}
    \setlength{\itemsep}{1em}
}{
    \end{enumerate}
}

\let\sectiononly\section
\renewcommand{\section}[1]{
    \sectiononly{#1}
    \let\sectiontitle\relax
    \newcommand{\sectiontitle}{#1}
}

\title{蓝色发光二极管}
\author{
    张植竣\texorpdfstring{\quad}{    }高乐耘\texorpdfstring{\quad}{    }付大为
}
\date{2022年4月29日}

\begin{document}

\captionsetup{font={small}}
\addtocounter{framenumber}{-1}
\begin{frame}[plain]
    \titlepage
\end{frame}

\date{}

\part{LED简介}

\begin{frame}
    \partpage
\end{frame}

\section{何为LED}

\begin{slide}
    A light-emitting diode (LED) is a semiconductor light source that emits light when current flows through it. Electrons in the semiconductor recombine with electron holes, releasing energy in the form of photons. The color of the light (corresponding to the energy of the photons) is determined by the energy required for electrons to cross the band gap of the semiconductor. White light is obtained by using multiple semiconductors or a layer of light-emitting phosphor on the semiconductor device.
    \bigskip

    发光二极管（LED）是一种\textbf{通电时发光的半导体光源}。半导体中的电子和空穴重新结合时以光子的形式释放出能量。光的颜色（与光子的能量对应）由电子穿过半导体的能带空隙所需的能量决定。白光由多种半导体同时发光或由一层发光磷光体在半导体设备上获得。
\end{slide}

\begin{slide}
    \begin{columns}
        \begin{column}{0.5\linewidth}
            \begin{figure}
                \centering
                \includegraphics[width=0.8\linewidth]{pictures/RBG-LED.pdf}
            \end{figure}
        \end{column}
        \begin{column}{0.5\linewidth}
            \begin{figure}
                \centering
                \includegraphics[width=0.8\linewidth]{pictures/LED,_5mm,_green_(en).pdf}
            \end{figure}
        \end{column}
    \end{columns}
    \bigskip
    \url{https://en.wikipedia.org/wiki/Light-emitting_diode} \\
    Leadframe（引线框架），die（芯片），Reflective cavity（反射腔），Wire bond（焊线），Epoxy case（环氧树脂外壳）
\end{slide}

\section{LED发展历史}

\begin{slide}
    Appearing as practical electronic components in 1962, the earliest LEDs emitted low-intensity infrared (IR) light. Infrared LEDs are used in remote-control circuits, such as those used with a wide variety of consumer electronics. The first visible-light LEDs were of low intensity and limited to red. Early LEDs were often used as indicator lamps, replacing small incandescent（发白热光的） bulbs, and in seven-segment displays.
    \bigskip

    最早的发光二极管于1962年作为实际使用的电子元件出现，可发出\textbf{低亮度的红外光}。红外LED被用于遥控电路中，例如制造各种各样的电子消费品。最早的可见光LED亮度低，且局限于红光。早期LED常用作台灯指示灯，替代原先的小白炽灯，以及制造七段数码显示管。
\end{slide}

\begin{slide}
    Recent developments have produced LEDs available in visible, ultraviolet (UV), and infrared wavelengths, with high, low, or intermediate light output, for instance white LEDs suitable for room and outdoor area lighting. LEDs have also given rise to new types of displays and sensors, while their high switching rates are useful in advanced communications technology with applications as diverse as aviation lighting, fairy lights, automotive headlamps, advertising, general lighting, traffic signals, camera flashes, lighted wallpaper, horticultural grow lights, and medical devices.
    \bigskip

    近期发展出多种LED，包括\textbf{可见光、紫外光、红外光波段，覆盖低中高各亮度}，对室内和室外空间均有适配。随LED还发展出多种显示和传感技术，其高转换速率在先进通信技术中大有用途，包括航空照明、小彩灯、汽车头灯、广告、一般照明、交通信号灯、相机闪光灯、发光墙纸、园艺灯和医学设备等。
\end{slide}

\section{LED的优劣}

\begin{slide}
    LEDs have many advantages over incandescent light sources, including lower power consumption, longer lifetime, improved physical robustness, smaller size, and faster switching. In exchange for these generally favorable attributes, disadvantages of LEDs include electrical limitations to low voltage and generally to DC (not AC) power, inability to provide steady illumination from a pulsing DC or an AC electrical supply source, and lesser maximum operating temperature and storage temperature.
    \bigskip

    相比白炽灯光源，LED有众多优势，比如更低功率损耗、更长寿命、更好的物理鲁棒性、更小的体积和更快的切换本领。作为补偿，LED的缺点包括低电压局限和直流电供电局限，在脉动直流电或交流电源下无法稳定发光，且最高工作和储存温度较低等。
\end{slide}

\begin{slide}
    In contrast to LEDs, incandescent lamps can be made to intrinsically run at virtually any supply voltage, can utilize either AC or DC current interchangeably, and will provide steady illumination when powered by AC or pulsing DC even at a frequency as low as 50 Hz. LEDs usually need electronic support components to function, while an incandescent bulb can and usually does operate directly from an unregulated DC or AC power source.
    \bigskip

    与LED相比，白炽灯固有地在几乎任意的供电电压下都可以工作，可以交替使用交流或直流电，在频率低至50赫兹的交流电或脉动直流电下仍能提供稳定照明。LED发光通常\textbf{需要其他电子元件的支持}，但白炽灯可以且通常直接接到非受控直流电源或交流电源上工作。
\end{slide}

\part{蓝光二极管}

\begin{frame}
    \partpage
\end{frame}

\section{2014年诺贝尔物理学奖}

\begin{slide}{\url{https://www.nobelprize.org/uploads/2018/06/press-25.pdf}}
    \begin{figure}
        \centering
        \includegraphics[width=0.8\textwidth]{pictures/Nobel.pdf}
    \end{figure}
    \bigskip

    The Royal Swedish Academy of Sciences（瑞典皇家科学院），\\
    Isamu Akasaki（赤崎勇，2021年4月1日逝世，享年92岁），\\
    Hiroshi Amano（天野浩），Shuji Nakamura（中村修二）
    \bigskip

    图及以下英文文字为诺贝尔奖官网对此次工作的简短介绍。
\end{slide}

\begin{slide}{\url{https://www.nobelprize.org/uploads/2018/06/ljussignal_whitesign.jpg}}
    \centering
    \includegraphics[width=.5\textwidth]{pictures/signal_whitesign.pdf}
\end{slide}

\begin{slide}{\url{https://www.nobelprize.org/uploads/2018/06/press-25.pdf}}
    When Isamu Akasaki, Hiroshi Amano and Shuji Nakamura produced \textbf{bright blue light beams from their semi-conductors} in the early 1990s, they triggered \textbf{a fundamental transformation of lighting technology}. Red and green diodes had been around for a long time but without blue light, white lamps could not be created. Despite considerable efforts, both in the scientific community and in industry, the blue LED had remained \textbf{a challenge for three decades}.
    \bigskip

    赤崎勇、天野浩和中村修二在20世纪90年代前期的工作不可不谓之巨大突破。他们使用半导体材料发出了\textbf{明亮的蓝光}，催生了\textbf{发光技术底层的转型}。那时红绿两色的发光二极管已有多年历史，唯有蓝色缺位，白光二极管一直无法制成。此前人们前赴后继\textbf{三十多年}，仍未能解决这个横跨科学界和工业界的\textbf{巨大难题}。
\end{slide}

\begin{slide}{\url{https://www.nobelprize.org/uploads/2018/06/press-25.pdf}}
    \textbf{They succeeded where everyone else had failed}. Akasaki worked together with Amano at the University of Nagoya, while Nakamura was employed at Nichia Chemicals, a small company in Tokushima. Their inventions were revolutionary. \textbf{Incandescent light bulbs lit the 20th century; the 21st century will be lit by LED lamps.}
    \bigskip

    \textbf{他们在万人跌倒之处成功了}！赤崎勇和天野浩一起在名古屋大学工作，而中村修二在日亚化学工业株式会社（日本四国岛东岸港市德岛的一家小公司）工作。他们的发明是革命性的。\textbf{白炽灯照亮了\textbf{20}世纪，而\textbf{21}世纪将为\textbf{LED}灯所点亮。}
\end{slide}

\begin{slide}{\url{https://www.nobelprize.org/uploads/2018/06/press-25.pdf}}
    White LED lamps emit a bright white light, are \textbf{long-lasting} and \textbf{energy-efficient}. They are constantly improved, getting more efficient with higher luminous flux (measured in lumen) per unit electrical input power (measured in watt). The most recent record is \textbf{just over} $\mathbf{300\,lm/W}$, which can be compared to 16 for regular light bulbs and close to 70 for fluorescent lamps.
    \bigskip

    白光二极管发光明亮，\textbf{持久}，且\textbf{节能}。它们不断随更新换代而变得越来越高效，即单位输入电功率（单位瓦特）获得的照度（单位流明）不断提高。其流明瓦特比\textbf{最近刚好超过了}$\mathbf{300\,lm/W}$，该指标对于常规灯泡而言仅16，荧光灯接近70。
\end{slide}

\begin{slide}{\url{https://www.nobelprize.org/uploads/2018/06/press-25.pdf}}
    \small
    As \textbf{about one fourth of world electricity consumption} is used for lighting purposes, the LEDs contribute to saving the Earth's resources. Materials consumption is also diminished as LEDs last \textbf{up to 100,000 hours}, compared to 1,000 for incandescent bulbs and 10,000 hours for fluorescent lights.
    \bigskip

    由于\textbf{世界上约四分之一的电都消耗在了照明用途上}，LED为全球节能带来了巨大贡献。由于LED寿命\textbf{长达十万小时}，相比之下白炽灯只有一千小时，荧光灯有一万小时，制造光源所需的材料损耗亦同步减少。
    \end{slide}

\begin{slide}{\url{https://www.nobelprize.org/uploads/2018/06/efficiency.pdf}}
    \begin{figure}[h]
        \centering
        \includegraphics[width=1\textwidth]{pictures/efficiency.pdf}
    \end{figure}
\end{slide}

\begin{slide}{\url{https://www.nobelprize.org/uploads/2018/06/press-25.pdf}}
    The LED lamp holds great promise for increasing the quality of life for \textbf{over 1.5 billion people} around the world who lack access to electricity grids: due to \textbf{low power requirements} it can be powered by \textbf{cheap local solar power}. The invention of the efficient blue LED is \textbf{just twenty years old}, but it has already contributed to create white light in an entirely new manner to the benefit of us all.
    \bigskip

    LED灯对提高全球\textbf{超过\textbf{1.5}亿与电网隔绝的人口}的生活质量大有前景：它\textbf{需要的电功率很小}，可以使用\textbf{当地便宜的太阳能}供电。高效蓝光LED的发明利在千秋，这项发明\textbf{不过只有二十年历史}，但已完全革新了人类产生白光的方式。
\end{slide}

\section{发展历史}

\begin{slide}{\url{https://en.wikipedia.org/wiki/Light-emitting_diode}}
    The first blue-violet LED using \textbf{magnesium-doped gallium nitride} was made at Stanford University in 1972 by Herb Maruska and Wally Rhines, doctoral students in materials science and engineering. At the time Maruska was on leave from RCA Laboratories, where he collaborated with Jacques Pankove on related work.
    \bigskip

    世界首个蓝紫光LED由美国斯坦福大学于1972年使用\textbf{掺镁氮化镓}制成，主要贡献者为Herb Maruska和Wally Rhines，他们是材料科学与工程专业的博士生。那时Maruska正休假，不在RCA实验室，在那个实验室里他与Jacques Pankove在发光二极管有关工作中合作。
\end{slide}

\begin{slide}{\url{https://en.wikipedia.org/wiki/Light-emitting_diode}}
    In 1971, the year after Maruska left for Stanford, his RCA colleagues Pankove and Ed Miller demonstrated \textbf{the first blue electroluminescence from zinc-doped gallium nitride}, though the subsequent device Pankove and Miller built, \textbf{the first actual gallium nitride light-emitting diode, emitted green light}.
    \bigskip

    1971年，也就是Maruska离开斯坦福大学的一年后，他的RCA同事Pankove和Ed Miller演示了\textbf{首次掺锌氮化镓的蓝色电致发光}，通过Pankove和Miller建造的后续装置，\textbf{首个货真价实的氮化镓LED发出了绿光}。
\end{slide}

\begin{slide}{\url{https://en.wikipedia.org/wiki/Light-emitting_diode}}
    In 1974 the U.S. Patent Office awarded Maruska, Rhines and Stanford professor David Stevenson a patent for their work in 1972. Today, \textbf{magnesium-doping of gallium nitride remains the basis for all commercial blue LEDs and laser diodes}. In the early 1970s, these devices were too dim for practical use, and research into gallium nitride devices slowed.
    \bigskip

    1974年，美国专利局给Maruska、Rhines和斯坦福教授David Stevenson一行人1972年的工作授予了专利。今天，\textbf{掺镁氮化镓仍然是商用蓝光二极管和激光二极管的基础}。20世纪70年代前期，这些设备的实用价值还不清楚，有关氮化镓设备的研究放缓。
\end{slide}

\begin{slide}{\url{https://en.wikipedia.org/wiki/Light-emitting_diode}}
    In August 1989, Cree introduced the first commercially available blue LED based on \textbf{the indirect bandgap semiconductor, silicon carbide (SiC)}. SiC LEDs had \textbf{very low efficiency}, no more than about 0.03\%, but \textbf{did emit in the blue portion of the visible light spectrum}.
    \bigskip

    1989年八月，Cree以\textbf{间接带隙半导体碳化硅}为基础引进了首个商用的蓝光二极管。碳化硅LED\textbf{效率极低}，不超过0.03\%，但它\textbf{真真切切地能在可见光波段发蓝光}。
\end{slide}

\begin{slide}{\url{https://en.wikipedia.org/wiki/Light-emitting_diode}}
    In the late 1980s, key breakthroughs in \textbf{GaN epitaxial growth and p-type doping} ushered in the modern era of GaN-based optoelectronic devices. Building upon this foundation, Theodore Moustakas at Boston University patented a method for producing high-brightness blue LEDs using a new two-step process in 1991.
    \bigskip

    二十世纪80年代后期，\textbf{氮化镓外延生长技术和p型掺杂技术}的关键突破引领了近代以氮化镓为基础的光电设备时代。在此基础上，1991年波士顿大学的Theodore Moustakas获得了使用新型两步反应产生高亮度蓝光二极管方法的专利。
\end{slide}

\begin{slide}{\url{https://en.wikipedia.org/wiki/Light-emitting_diode}}
    Two years later, in 1993, \textbf{high-brightness blue LEDs were demonstrated by Shuji Nakamura of Nichia Corporation using a gallium nitride growth process}. In parallel, Isamu Akasaki and Hiroshi Amano of Nagoya University were working on \textbf{developing the important GaN deposition on sapphire substrates and the demonstration of p-type doping of GaN}.
    \bigskip

    两年后，也就是1993年，\textbf{高亮蓝光二极管由日本日亚化学工业株式会社的中村修二使用氮化镓生长过程制成}。与此同时，名古屋大学的赤崎勇和天野浩正致力于\textbf{研发氮化镓在蓝宝石衬底上的重要沉积和氮化镓p型掺杂的实现}。
\end{slide}

\begin{slide}{\url{https://en.wikipedia.org/wiki/Light-emitting_diode}}
    This new development revolutionized LED lighting, making high-power blue light sources practical, leading to the development of technologies like Blu-ray.
    \bigskip

    这些新进展对LED发光技术带来了革命性影响，使得制造高功率蓝光源变得切实可行，引领了蓝光光碟等技术的发展。
\end{slide}

\begin{slide}{\url{https://en.wikipedia.org/wiki/Light-emitting_diode}}
    Nakamura was awarded the 2006 Millennium Technology Prize for his invention. Nakamura, Hiroshi Amano and Isamu Akasaki were awarded the Nobel Prize in Physics in 2014 for the invention of the blue LED.
    \bigskip

    中村修二因其发明被授予2006年度的千禧年科技奖。中村修二、赤崎勇和天野浩三人因蓝光二极管的发明而被授予2014年度的诺贝尔物理学奖。
\end{slide}

\section{发光原理和技术难点}

\begin{slide}{\url{https://www.nobelprize.org/uploads/2018/06/diode.pdf}}
    \includegraphics[width=\textwidth]{pictures/diode.pdf}
    Sapphire Substrate: 蓝宝石基片
\end{slide}

\begin{slide}{\url{https://www.azom.com/article.aspx?ArticleID=11451}}
    On 7th October 2014, Isamu Akasaki, Hiroshi Amano and Shuji Nakamura were awarded the Nobel Prize in Physics for the invention of efficient blue light-emitting diodes (LEDs).
    \medskip
    
    Red and green LEDs were created in a number of laboratories during the 1950s and 1960s but it took \textbf{another three decades to finally produce efficient blue LEDs} -- why were they so hard to make?
    \bigskip

    以下这篇文章介绍了蓝光二极管的结构、发光原理及两项主要的技术困难。其中也穿插了不少科学史。
    \medskip

    蓝光二极管比红光和绿光二极管的到来晚了三十年，这一诺贝尔奖的工作为何如此困难？
\end{slide}

\begin{slide}{\url{https://www.azom.com/article.aspx?ArticleID=11451}}
    Light-emitting diodes are electronic devices which are \textbf{illuminated by the movement of electrons in a semiconductor material}. LEDs are able to emit light with a range of wavelengths from the infrared to the ultraviolet. A semiconductor is a material with an electrical conductivity which is \textbf{somewhere between a conductor such as copper}, and an \textbf{insulator such as rubber}. They are usually \textbf{made from a poor conductor which is then ‘doped’ by adding atoms of another material to it}.
    \bigskip

    简单地说，LED发光源于半导体材料中电子的移动（电致发光）。所谓半导体，其导电能力介于导体和绝缘体之间。半导体常由不良导体掺杂制成。所谓掺杂，就是在材料中掺入本不属于该材料组分的其他原子。
\end{slide}

\begin{slide}{\url{https://www.azom.com/article.aspx?ArticleID=11451}}
    LEDs are typically made from \textbf{aluminum-gallium-arsenide (AlGaAs)} which in its pure form does not contain any free electrons to conduct electrical current. As a result, AlGaAs is \textbf{doped with either free electrons or ‘electron holes’ in order to change the materials balance and make it more conductive}.
    \medskip

    What is an \textbf{Electron Hole}? The absence of an electron from an \textbf{otherwise full valence band}.
    \bigskip

    一般LED用AlGaAs制成，这种材料本身不带有可以导电的自由电子。人们在其中掺杂，使其内部含有自由电子或电子空穴，改变其物料平衡，改善其导电性。所谓电子空穴，即因中心原子的共价键缺失导致的缺电子现象，是一种物理模型，而非真实存在的物理实体。
\end{slide}

\begin{slide}{\url{https://www.azom.com/article.aspx?ArticleID=11451}}
    Semiconductors can be classified into two types of material; N-type and P-type:
    \medskip

    \textbf{N-type} semiconductors contain extra negatively charged \textbf{electrons} and as a result the free electrons flow from negatively charged areas to positively charged areas.
    \medskip

    \textbf{P-type} semiconductors have \textbf{extra holes which allows free electrons to jump between the holes} and moving from negatively charged areas to positively charged areas as a result.
    \bigskip

    因掺杂类型的不同，半导体分为N型（negtive）和P型（positive）。N型半导体的载流子为富余电子，P型半导体的载流子为富余空穴。空穴导电并非空穴本身移动，而是电子在不同空穴间跳转，等效地看，就是空穴在向电子运动的反方向移动。
\end{slide}

\begin{slide}{\url{https://www.azom.com/article.aspx?ArticleID=11451}}
    A diode consists of a section of an N-type semiconductor attached to a section of a P-type semiconductor (known as a \textbf{p-n junction}) with two electrodes (电极) placed at either end of this arrangement. When current flows across a diode, the negatively charged electrons move in one direction in the material and the positively charged holes move in the opposite direction. As the holes exist in lower energy states, \textbf{a free electron will lose energy when it falls to a hole and emit this energy in the form of a photon of light}.
    \bigskip

    二极管由P-N结组成。所谓P-N结，即两块结合在一起的不同掺杂类型的半导体。电流通过P-N结时，电子和空穴分别向相反的方向移动。电子与空穴可能发生结合，该过程产生能量，并可能以光子的形式放出，这样便会发光。
\end{slide}

\begin{slide}{\url{https://www.azom.com/article.aspx?ArticleID=11451}}
    \textbf{The size of the fall in energy determines the energy the photon has when it is emitted, which in turns determines the colour of the light the diode emits}. An emitted photon with a large amount of energy will have a shorter wavelength than light emitted with a lower amount of energy.
    \bigskip

    电子空穴结合能决定了结合过程释放的能量，进而决定了可能释放出的光子的能量，从而影响发光的颜色。
\end{slide}

\begin{slide}{\url{https://www.azom.com/article.aspx?ArticleID=11451}}
    \textbf{The first diode} which was able to emit \textbf{electrically produced light} was created in 1907 by H.J. Round whilst he was experimenting with a cat's-whisker detector. Round applied a potential difference across a \textbf{silicon carbide (SiC) crystal}. He found that \textbf{the colour of light emitted varied depending on the voltage which was applied across the crystal}.
    \bigskip

    碳化硅是一种可用于生产半导体发光二极管的材料。早在1907年，这种材料就被用于制造最早的LED。碳化硅晶体的发光颜色与供电电压有关，也可发出蓝光，而现在之所以不用它，如前所述，其发光效率太低了。
\end{slide}

\begin{slide}{\url{https://www.azom.com/article.aspx?ArticleID=11451}}
Electroluminescence was widely studied in detail during the 1950s by scientists, with J.R. Haynes demonstrating that \textbf{the emission of light from germanium and silicon diodes was due to the interaction of holes and electrons in a p-n junction} in 1956 at Bell Telephone Laboratories. In 1956, Shockley, Bardeen and Brattain received a Nobel Prize for their research into semiconductors and the discovery of the transistor effect.
\bigskip

1950年代，科学家对电致发光进行了广泛的详细研究，Haynes于1956年在贝尔电话实验室证明了锗和硅二极管的光发射是由于p-n结中空穴和电子的相互作用。1956年，Shockley、Bardeen和Brattain因对半导体的研究和晶体管效应的发现而获得诺贝尔奖。

\end{slide}

\begin{slide}{\url{https://www.azom.com/article.aspx?ArticleID=11451}}
    \textbf{Infrared LEDs were created in 1962 using p-n junctions made from GaAs}. This semiconductor material has a direct bandgap of \textbf{1.4 eV}, which directly corresponds to the wavelength of infrared light. By the end of the 1960s, \textbf{red and green LEDs were being manufactured in different countries using p-n junctions made from GaP}. The development of a blue LED however proved far more difficult to scientists.
    \bigskip

    另一种半导体材料砷化镓被用于制造红外LED，其直接带隙为1.4 eV，落在红外区。还有一种半导体材料磷化镓，适当掺杂做成二极管可以使其发出红光和绿光。而对蓝光而言，就没有那么容易了。
\end{slide}

\begin{slide}{\url{https://www.azom.com/article.aspx?ArticleID=11451}}
    \begin{multicols}{2}
    \scriptsize
    \begin{figure}
        \centering
        \includegraphics[width=0.6\columnwidth]{pictures/Wurtzite-unit-cell-3D-balls.pdf}
    \end{figure}
    \bigskip
    
    The first attempts at \textbf{the emission of blue light from a diode used ZnSe and SiC}, which are semiconductors with \textbf{high indirect bandgaps, but did not produce efficient light emission}. The material which enabled the development of blue LEDs was \textbf{gallium nitride (GaN)}. GaN is a semiconductor with a Wurtzite（纤锌矿） crystal structure and a \textbf{direct bandgap of 3.4 eV}, which directly corresponds to the wavelength of light in the \textbf{ultraviolet} range.
    \bigskip

    为了使二极管发蓝光，我们需要寻找高带隙半导体。最先应用的硒化锌和碳化硅制造出的LED发光效率太低。高效的蓝光LED始于氮化镓，这是一种纤锌矿型半导体，其直接带隙3.4 eV，在紫外区。
    \end{multicols}
    {\tiny 图片来源：\url{https://en.wikipedia.org/wiki/File:Wurtzite-unit-cell-3D-balls.png}
}
\end{slide}

\begin{slide}{\url{https://www.azom.com/article.aspx?ArticleID=11451}}
    \textbf{At the end of the 1950s}, researchers at Philips Research Laboratories had considered using GaN to create blue LEDs but instead chose to \textbf{concentrate on devices made from GaP instead due to difficulties in growing GaN crystals}. These crystals were \textbf{more effectively produced in the late 1960s by using the Hydride Vapour Phase Epitaxy （HVPE, 氢化物气相外延） growing GaN on a substrate}.
    \bigskip

    早在20世纪50年代末，氮化镓便被考虑用于制作蓝光LED，但由于其晶体生长过于困难，飞利浦实验室选择了磷化镓。直至20世纪60年代后期，氮化镓得以高效生产——利用氢化物气相外延法在基片上生长氮化镓晶体。
\end{slide}

\begin{slide}{\url{https://www.azom.com/article.aspx?ArticleID=11451}}
    In 1974 Isamu Akasaki began studying gallium nitride and took up a professorship at Nagoya University to continue his research alongside Hiroshi Amano. In 1986 the \textbf{MOVPE（金属有机物气相外延）} technique was used in order to produce GaN with high crystal quality and good optical properties. Shuji Nakamura later development a similar method in order to grown GaN at low temperatures.
    \bigskip

    可以推断，虽然氮化镓晶体生长在工业上已经实现，其技术水平尚不满足蓝光二极管的制作需求（不然为什么又折腾了三十年）。1974年开始，2014年诺贝尔物理学奖得主们赤崎勇和天野浩开始研究之。1986年，\textbf{金属有机物气相外延法}的出现大大改善了该半导体材料的光学特性。另一位得奖者中村修二随后以低温方法实现了相似技术（该技术对于生产蓝光二极管更为高效）。
\end{slide}

\begin{slide}{\url{https://www.azom.com/article.aspx?ArticleID=11451}}
    Another major problem in producing blue LEDs was \textbf{the difficulty in p-doping GaN with precision}. In the late 1980s, Amano and Akasaki discovered that \textbf{when GaN was doped with zinc atoms, it emitted more light} and thus this gave better p-doping. This phenomena was later explained in an article by Nakamura. This was an important discovery as it \textbf{paved the way for p-n junctions to be used in gallium nitride semiconductors}.
    \bigskip

    除氮化镓晶体生长技术外，蓝光二极管生产的另一技术难题是精确p型掺杂问题。该掺杂技术于20世纪80年代后期取得突破，氮化镓掺锌会发更多光，从而更有利于p型掺杂\textbf{（为什么？）}。这一突破为氮化镓材料在蓝光二极管的应用奠定了基础。
\end{slide}

\begin{slide}{\url{https://www.azom.com/article.aspx?ArticleID=11451}}
    A \textbf{key step} in the development of blue LEDs was the development of heterojunctions（异质结） in the early 1990s by research groups led by Akasaki and Nakamura. \textbf{In 1994, Nakamura used a double heterojunction InGaN/AlGaN to produce a device with a quantum efficiency of 2.7\%, which opened the door for efficient blue LEDs to be easily produced}.
    \bigskip

    蓝光二极管发展成熟的关键一步是异质结的应用。这一重大突破产生于20世纪90年代前期，由二位诺贝尔奖得主赤崎勇和中村修二所贡献。1994年，中村修二使用了双异质结InGaN/AlGaN来制作蓝光LED，其量子效率\textbf{（是什么意思？）}达到了2.7\%，开创了蓝光LED量产的新时代。
\end{slide}

\begin{slide}{\url{https://www.azom.com/article.aspx?ArticleID=11451}}
    The emission of blue light based on GaN was \textbf{observed between 1995 and 1996 by both research groups}. Modern efficient GaN based LEDs are the result of long series of breakthroughs in the scientific fields of materials physics, optics, electronics and chemistry.
    \bigskip

    终于，实验者们在1995年和1996年分别使氮化镓二极管发出了蓝光！如今蓝光二极管的高效离不开长期材料物理、光学、电子学和化学领域的一系列重大突破。
\end{slide}

\begin{slide}{\url{https://www.chemeurope.com/en/news/151134/}}
    Scientists at UCL, in collaboration with groups at the University of Bath and the Daresbury Laboratory, have uncovered the mystery of why blue light-emitting diodes (LEDs) are so difficult to make, by revealing \textbf{the complex properties of their main component - gallium nitride} - using sophisticated computer simulations.
    \bigskip

    下面通过一篇PRL文献的工作介绍，展示生产掺杂氮化镓晶体材料的困难。
    \medskip
    
    该文献利用复杂的计算机模拟方法，描绘了向氮化镓晶体中掺入镁原子时晶体发生的复杂变化，揭示了掺杂技术的难点所在。
\end{slide}

\begin{slide}{\url{https://www.chemeurope.com/en/news/151134/}}
    The key ingredient for blue LEDs is \textbf{gallium nitride, a robust material with a large energy separation, or `gap', between electrons and holes - this gap is crucial in tuning the energy of the emitted photons to produce blue light}. But while doping to donate mobile negative charges in the substance proved to be easy, \textbf{donating positive charges failed completely}. The breakthrough, which won the Nobel Prize, \textbf{required doping it with surprisingly large amounts of magnesium}.
    \bigskip

    刚才我们提到的获得诺贝尔奖的工作，即双异质半导体蓝光二极管的发明，离不开掺杂技术的突破。刚才我们也提到，p型掺杂是一大技术难题。我们需要在氮化镓晶体中大量地掺入镁原子，让我们来看看这是为什么。
\end{slide}

\begin{slide}{\url{https://www.chemeurope.com/en/news/151134/}}
    But in fact, \textbf{to donate a single mobile hole in gallium nitride, at least a hundred atoms of magnesium have to be added}. It's technically extremely difficult to manufacture gallium nitride crystals with so much magnesium in them, not to mention that it's been frustrating for scientists not to understand what the problem was.
    \bigskip

    出乎意料地，模拟实验表明，为了在氮化镓晶体中产生一个正电荷，需要掺入上百的镁原子。掺杂如此之多，在工业上是个困难。然而，其缘由无人能解。
\end{slide}

\begin{slide}
    \begin{multicols}{2}
    \parbox{\columnwidth}{
        \includegraphics[width=\columnwidth]{pictures/cat.pdf} \\
        \tiny 图片来源: \url{https://pkuhelper.pku.edu.cn/hole}
    }
    \kaishu
    \raggedright\qquad 我可以有把握地说，没人懂得量子力学。 \\
    \raggedleft ——理查德·费曼 \\
    ~\\
    \raggedright\qquad 量子力学绝对是上个世纪最伟大的发现，没有之一。没人能弄懂量子力学，但所有这些半导体元件都是量子力学创造出来的。看看，我们的手机，我们的电脑，里面有多少三极管。 \\
    \raggedleft ——赵建业（学生回忆版） \\
    \end{multicols}
\end{slide}

\begin{slide}{\url{https://www.chemeurope.com/en/news/151134/}}
    Such models have been widely applied to the study of perfect crystals, where a small group of atoms form a repeating pattern. \textbf{Introducing a defect that breaks the pattern presents a conundrum （难题）, which required the UK's largest supercomputer to solve.} Indeed, \textbf{calculations on very large numbers of atoms were therefore necessary but would be prohibitively expensive to treat the system on a purely quantum-mechanical level}.
    \bigskip

    如何进行模拟计算呢？纯量子模型本适用于理想晶体，然而掺杂的引入破坏了系统的对称性，如果使用纯量子力学模型，计算量在英国顶级超算级别。实际上，在计算机科学中，这几乎是一个不可解问题，太多的原子需要参与计算了。
\end{slide}

\begin{slide}{\url{https://www.chemeurope.com/en/news/151134/}}
    The team's solution was to apply an approach pioneered in another piece of Nobel Prize winning research: \textbf{hybrid quantum and molecular modelling, the subject of 2013's Nobel Prize in Chemistry}. In these models, different parts of a complex chemical system are simulated with different levels of theory.
    \bigskip

    必须更换模型！团队使用了获得2013年诺贝尔化学奖的主要工作——混合量子分子模型，该模型使得复杂化学系统的计算模拟成为可解问题。该模型使用不同层次的理论模拟复杂化学系统的不同部分，从而大大降低了计算量。
\end{slide}

\begin{slide}{\url{https://www.chemeurope.com/en/news/151134/}}
    The simulation tells us that when you \textbf{add a magnesium atom}, it replaces a gallium atom but \textbf{does not donate the positive charge to the material}, instead keeping it to itself,' says Richard Catlow (UCL Chemistry), one of the study's co-authors. `In fact, to provide enough energy to \textbf{release the charge} will require heating the material \textbf{beyond its melting point}.
    \bigskip

    每掺入一个镁原子，将替换掉晶体中的一个镓管子，但可能并不会贡献一个正电荷，相反地，镁原子自己保留了正电荷。想要释放这一正电荷，靠热运动是不可能的，因为电离能的大小超过了熔点下的热能。
\end{slide}

\begin{slide}{\url{https://www.chemeurope.com/en/news/151134/}}
    Even if it were released, it \textbf{would knock an atom of nitrogen out of the crystal, and get trapped anyway in the resulting vacancy}. Our simulation shows that the behaviour of the semiconductor is \textbf{much more complex than previously imagined, and finally explains why we need so much magnesium to make blue LEDs successfully}.
    \bigskip

    即使正电荷释放出来，也可能经历一定的过程挤出一个氮原子并将电荷束缚于氮原子留下的空隙中。模拟表明，半导体行为的复杂程度远超我们想象，镁原子的如此“浪费”解释了需要大量掺杂的原因。
\end{slide}

\begin{slide}{\url{https://www.chemeurope.com/en/news/151134/}}
    The simulations crucially fit a complete set of previously unexplained experimental results involving the behaviour of gallium nitride. Aron Walsh (Bath Chemistry) says `We are now \textbf{looking forward to the investigations into heavily defective GaN, and alternative doping strategies to improve the efficiency of solid-state lighting.}
    \bigskip

    模拟结果很好地解释了氮化镓掺杂实验中观察到的镁原子“浪费”现象，下一步应当深入研究重掺杂氮化镓以及其他备选的掺杂策略以改进固态照明效率。
\end{slide}

\begin{slide}{技术难点小结}
    \begin{enumerate}
    \item 难题之零为找到合适的半导体材料——氮化镓等
    \item 难题之一为特定理化性质的氮化镓等晶体层的制备
    \item 难题之二为氮化镓等半导体材料n型重掺杂的工业实现
    \end{enumerate}
    \bigskip

    \begin{enumerate}
    \item 难题之零随时代发展和技术进步迎刃而解，当今氮化镓等材料的广泛应用是难题一二求解的必然结果
    \item 难题之一困扰人类三十年之久，经过几代科学家和技术人员的努力，最终由中村修二等人解决
    \item 难题之二与难题之一相伴而生，在三人掺锌和使用双异质结等新思路下，n型掺杂效率问题得以解决
    \end{enumerate}
\end{slide}

\part{从蓝光二极管到OLED的现代化进程}

\begin{frame}
    \partpage
\end{frame}

\section{从蓝光LED到白光LED}

\begin{slide}
\begin{enumerate}[1.]
    \item 1961年美国公司德州仪器的Robert Biard与Gary Pittman首次发现了砷化镓及其他半导体合金的红外放射作用,也就是红光LED,标志着LED的诞生
    \item 过了30多年, 到1993年蓝光LED的发明, 才标志着LED的全彩时代拉开帷幕
    \item 蓝光二极管的重要性不仅在于它本身,更在于推动了白光二极管的诞生,也就是目前广泛用于生产生活的灯源
\end{enumerate}
\end{slide}

\section{白光LED的制造}

\begin{slide}
有了能发出蓝光的发光二极管就可以制造白光LED，主要有以下两种制造方式:
\begin{enumerate}[1]
    \item 这种方法涉及用不同颜色的荧光粉涂覆一种颜色的LED（主要是由InGaN制成的蓝色LED）以形成白光；由此产生的LED称为基于荧光粉或荧光粉转换的白光LED（pcLED）。一小部分蓝光经历Stokes位移，将其从较短的波长转变为更长的波长。使用几个不同颜色的荧光粉层可以扩大发射光谱，有效地提高显色指数（CRI）。\\
    \textit{优点：制造简单，亮度高}\\
    \textit{缺点：LED发出的光能有一部分被转化成热能，效率较低}
    
    \item 结合蓝光发光二极管、红光发光二极管和绿光发光二极管便可做出白光发光二极管\\
    \textit{优点:有较广的色域,并且发光效率较高}\\
    \textit{缺点:具有较高的制作成本}
\end{enumerate}
\end{slide}

\begin{slide}{荧光粉基发光二极管}
    \small
    第一种白光LED采用了掺杂铈的钇铝石榴石（Ce:YAG）作为高亮度蓝色InGaN二极管上的涂层，这是一种荧光粉，在受到蓝色光照射时会发出光谱较宽的黄光（Stokes偏移）。它可以将部分蓝光转换为黄色，然后一起显示为白色。
    \begin{figure}
        \centering
        \includegraphics[width=0.5\linewidth]{pictures/700px-White_LED.pdf}
        \caption{白光LED的光谱显示基于GaN的LED直接发射的蓝光（峰值约为$465\,\mathrm{nm}$）和Ce:YAG荧光粉发出的更宽带的Stokes偏移光，其发射波长约为$500-700\,\mathrm{nm}$}
    \end{figure}
\end{slide}

\begin{slide}{荧光粉基发光二极管}
    基于荧光粉的LED由于Stokes位移和其他与磷光体相关的问题导致的热损失而产生效率损失。与普通LED相比，它们的发光效率取决于所产生光输出的光谱分布和LED的原始波长。截至2010年，最有效的黄色荧光粉仍然是YAG荧光粉，Stokes位移损耗不到10\%。
    \bigskip

    由于制造简单，荧光粉法仍然是制造高强度白光LED的最常用方法。使用具有荧光粉转换的单色发射器设计和生产光源比复杂的RGB系统更简单，更便宜，目前市场上的大多数高强度白光LED都是使用荧光粉光转换制造的。
\end{slide}

\begin{slide}{荧光粉基发光二极管}
    还有一种荧光粉基发光二极管类似于日光灯，发光单元是近紫外发光二极管，外面包着两种磷光剂混合物，一种是发红光和蓝光的铕，另一种磷光剂是发绿光的铜和铝掺杂了硫化锌。内里的紫外光发光二极管发出的紫外光被外层的磷光剂转换成红、蓝、绿三色光，混合后就成了白光。
    \bigskip

    \textit{优点:光谱的特性较佳，产生的光比较好看}\\
    \textit{缺点:紫外线会使黏合剂中的环氧树脂劣化变质，所以生产难度较高，而寿命亦较短,光波在转化过程中有较多被化成热能，因此效率较低}
\end{slide}

\begin{slide}{RGB系统}
    混合红色、绿色和蓝色光源以产生白光需要电子电路来控制颜色的混合。由于不同颜色发光模式略有不同，即使RGB光源封装在单个LED中，色彩平衡也可能根据视角而改变，因此RGB二极管很少用于产生白光。尽管如此，由于混合不同颜色的灵活性，这种方法有许多应用，原则上，这种机制在产生白光方面也具有更高的量子效率。
    \bigskip

    有几种类型的多色白光LED：二色、三色和四色白光LED。在这些不同方法中起作用的几个关键因素包括颜色稳定性，显色能力和发光效率。通常，更高的效率意味着更低的显色性，从而在发光效率和显色性之间进行了权衡。其中一个挑战是开发更高效的绿色LED，它的理论最大值为683 lm/W，但很少有绿色LED超过每瓦100流明，蓝色和红色LED接近其理论极限。
\end{slide}

\begin{slide}{RGB系统}
    \begin{columns}
        \begin{column}{0.5\linewidth}
            \begin{figure}
            \centering
            \includegraphics[width=0.8\linewidth]{pictures/440px-RGB_LED.pdf}
            \caption{封装在一起的红色和绿色LED}
            \end{figure}
        \end{column}
        \begin{column}{0.5\linewidth}
            \begin{figure}
            \centering
            \includegraphics[width=0.8\linewidth]{pictures/440px-RGB_LED_Spectrum.pdf}
            \caption{蓝色、黄绿色和高亮度红色固态半导体 LED 的组合光谱曲线。光谱带宽约为24 -- 27 nm}
            \end{figure}
        \end{column}
    \end{columns}
\end{slide}

\begin{slide}{RGB系统}
    多色LED还提供了一种形成不同颜色的光的新方法。大多数可感知的颜色可以通过混合不同量的三原色来形成。这允许精确的动态色彩控制。然而，这种类型的LED的发射功率随着温度的升高而呈指数级衰减，导致颜色稳定性发生巨大变化，这抑制了它的工业用途。
    \bigskip

    没有荧光粉的多色LED无法提供良好的显色性，因为每个LED都是窄带光源。没有荧光粉的LED虽然是一种较差的普通照明解决方案，但对于基于LED直接像素的显示器来说，这是最佳解决方案。
\end{slide}

\section{发光二极管(LED)的应用}

\begin{slide}
\textbf{状态显示灯}
    \textit{发光二极管所需推动电压及功率低，方便由运作电压低的微处理器控制及在以电池作电源的设备上使用，所以常被用在各种电子产品、设备的状态指示灯。在消费性电子产品，手提嵌入式电子设备，家庭电器、玩具、各种仪器……等用途上作为工作状态显示灯。}
    \begin{figure}[h]
        \centering
        \includegraphics[height=3cm, width=4.5cm]{pictures/Seven_segment_02_Pengo.pdf}
        \caption{用来制作数字显示器的发光二极管元件}
    \end{figure}
\end{slide}

\begin{slide}
\textbf{交通道路指示灯}
\begin{enumerate}
    \item 在车辆上的应用：\\
    \textit{当LED的光度发展达至要求后，便有汽车用LED作转向与刹车灯，主要好处是LED极高的开关速度，其亮起时间比白炽灯快$\mathrm{0.5s}$之多，且LED灯在白天时与白炽灯炮相比亮度较高、使用寿命长、点亮速度快，更可增加$\mathrm{6m}$的刹车辨识距离，也可降低事故发生，这对行车安全非常重要。而车辆行驶环境中震荡、温差变化下，LED仍然有稳定的光度及可靠性。}
\end{enumerate}
\end{slide}

\begin{slide}
    \textbf{交通道路指示灯}
    \begin{enumerate}
        \item 在道路上的应用：\\
        \textit{同样工作在户外环境的交通灯使用开关快速的LED也很适合，LED寿命较长也减少坏灯及影响交通的机会。}
        \item 公共场所的平板显示器：\\
        \textit{在机场、机舱、火车站、巴士站、码头等…各种公共交通工具上等地方，LED被普遍地采用作为平板显示器以显示如班次、目的地、时间等相关讯息。而LED的可靠性与低耗电，使其也适合用作紧急出口指示灯。}
    \end{enumerate}
    \end{slide}

\begin{slide}
\textbf{LED照明}\\
由于发光二极管灯被称为固态照明。传统照明灯俱如萤光灯、白炽灯及卤素灯都有装载气体的脆弱玻璃管，因而都不及全固态的LED坚固耐用。由于发光二极管在增加光度时，效率会下降，所以有些LED灯使用多个白光LED组合成一簇构成一个光源，以增加效率；同样的原因，在照明方面只用在对光度要求低的地方，以保持其较佳效率（省电）的特性，这些用途包括：
\begin{enumerate}
    \item 手电筒
    \item 紧急照明
    \item 手机闪光灯
    \item 道路及室内照明
\end{enumerate}
\end{slide}

\begin{slide}
\textbf{显示器}
\begin{enumerate}
    \item \textit{市面上现存普及的LED背光液晶显示电视的LED只用作背光光源，严格上并不是LED显示器，细小高解像度的视讯LED显示器现都采用OLED，OLED具有自发光性、广视角、高对比、低耗电、高反应速率等优点，OLED显示器因为不需背光源，所以可以比LCD显示器造得更薄}
    \item \textit{大型的LED显示器已普及于户外户内，户外LED显示器对解析度要求较低，但需要较高的亮度，多采用分立单色的LED组成。}
\end{enumerate}
\end{slide}

\section{有机发光二极管(OLED)}
\begin{slide}
有机发光二极体OLED。其发光原理跟发光二极管一样，不同之处是其发光物半导体是有机化合物（有机半导体），例如有机聚合物等。OLED制造流程简单，成本较低，可以用印刷等廉价生产方法制造，其优点包括：
\begin{enumerate}
    \item \textit{可以制造出大面积的面光源}
    \item \textit{元件本身可以是柔软、透明的}
\end{enumerate}
这些特性都是一般二极管所不及的。因此OLED可以造出大面积的照明灯具，软身、透明的显示器。
\end{slide}

\begin{slide}
有机发光二极体基本结构是由一薄而透明具半导体特性之铟锡氧化物（ITO），与电力之正极相连，再加上另一个金属阴极，包成如三明治的结构。整个结构层中包括了：电洞传输层（HTL）、发光层（EL）与电子传输层（ETL）。当电力供应至适当电压时，阳极空穴与阴极电子便会在发光层中结合，产生光子，依其材料特性不同，产生红、绿、蓝三原色，构成基本色彩。
    \begin{figure}[h]
        \centering
        \includegraphics[height=2cm, width=6cm]{pictures/OLED_schematic.pdf}
        \caption{OLED基本结构：1. 阴极 (−)；2. 发光层（Emissive Layer, EL）；3. 阳极空穴与阴极电子在发光层中结合，产生光子；4. 导电层（Conductive Layer）；5. 阳极 (+)}
    \end{figure}
\end{slide}

\begin{slide}
最早的OLED技术研发开始于1950年代的法国南茜大学，法国物化学家安德烈·贝纳诺斯获誉为“OLED之父”，最早的实用性OLED于1987被柯达公司的中国香港人邓青云和美国人Steven Van Slyke两人发现。

OLED英文名为Organic Light-Emitting Diode，缩写：OLED，中文名“有机发光二极管”都是邓青云命名的。
\medskip

\small\textit{邓青云自1975年开始加入柯达公司Rochester实验室从事有机发光二极体的研究工作，在意外中发现有机发光二极体。1979年的一天晚上，他在回家的路上忽然想起有东西忘记在实验室，回到实验室后，他发现在黑暗中的一块做实验用的有机蓄电池在闪闪发光，从而开始了对有机发光二极体的研究。1987年，邓青云和同事Slyke成功地使用类似半导体PN结的双层有机结构第一次作出了低电压、高效率的光发射器。为柯达公司生产有机发光二极体显示器奠定了基础。}
\end{slide}

\begin{slide}
OLED的特性是自发光，不像薄膜电晶体液晶显示器需要背光，因此可视度和亮度均高，且无视角问题，其次是驱动电压低且省电效率高，再加上反应快、重量轻、厚度薄，构造简单，成本低等特点，被视为21世纪最具前途的产品之一。
\bigskip

OLED可应用于制造平价可弯曲显示器、照明设备、发光衣或装饰墙壁。2004年开始，有机发光二极体已广泛应用于随身MP3播放器。不同颜色的OLED有不同寿命，衰退程度也不同（蓝色OLED的寿命最短），因此作为全彩色显视器时，色温会随使用时间而变，较常用的像点会较其他像点衰退得较快而使得光暗不均和偏色，也就是OLED显示器的“烧屏”（screen burn-in）现象。水份、湿气等会对OLED造成破坏，因此对封装的气密性、防水性也有要求。
\end{slide}

\begin{slide}
从技术上讲，问题在于蓝色LED的发光效率明显低于红色或绿色像素。这意味着，蓝色LED需要以更高的电流驱动，才能达到与红色或绿色像素相同的亮度。更高的电流会导致像素更快地退化，缩短其寿命，从而最终使显示屏的颜色向红色和绿色倾斜。因此OLED显示屏的颜色并不是均匀退化的，它最终会以红/绿色调呈现。因此，如果面板的某一区域长期显示蓝色或白色的图像，那么这个区域的蓝色像素会比其他区域退化得更快。这是烧屏产生的根本原因。
\bigskip

一般来讲，蓝色的OLED寿命最短，绿色OLED的寿命最长，因此厂家通过使蓝色、红色子像素变大，这样就需要更少的电流来驱动以提供所需的光，以延长它们的寿命，使得需要更长的时间才会出现明显的色彩变化。屏幕厂商据此研发出了Pentile排列和Diamond Pentile排列。
\end{slide}

\begin{slide}
    \begin{columns}
        \begin{column}{0.5\linewidth}
            \begin{figure}
            \centering
            \includegraphics[width=0.7\linewidth]{pictures/screen-burn-in-phone.pdf}
            
\includegraphics[width=0.7\linewidth]{pictures/Screen_Burn-in.pdf}
            \caption{手机和电视屏幕的烧屏现象}
            \end{figure}
        \end{column}
        \begin{column}{0.5\linewidth}
            \begin{figure}
            \centering
            \includegraphics[width=0.7\linewidth]{pictures/Pentile.pdf}
            \includegraphics[width=0.7\linewidth]{pictures/Diamond_Pentile.pdf}
            \caption{Pentile排列（上）和Diamond Pentile排列（下）}
            \end{figure}
        \end{column}
    \end{columns}
\end{slide}

\begin{slide}
此外，OLED还在光医学领域有广泛前景:
\begin{enumerate}
    \item 治疗创伤（PBM光疗法）：相比传统光医疗，OLED 贴片凭借二维弯曲（柔性）、轻薄、可穿戴等特点，可以更贴合皮肤，舒适感和体验性会更佳。
    \item 可穿戴OLED发光器件在个人移动医疗监测领域来测量人体信号，例如监测心跳和血氧水平。这是因为在心血管监测方法中，光体积描记法（PPG）信号和血氧饱和度（SpO2）水平是通过使用发光器件和光检测器，进行无创光测量的。
\end{enumerate}
\end{slide}


\addtocounter{framenumber}{-1}
\begin{frame}[plain]{参考文献}
\begin{thebibliography}{99}\setlength{\itemsep}{0pt}\small
    \bibitem{bib:LED} \url{https://en.wikipedia.org/wiki/Light-emitting_diode}
    \bibitem{bib:Nobel}
    \url{https://www.nobelprize.org/uploads/2018/06/press-25.pdf}
    \bibitem{bib:OLED}
    \url{https://en.wikipedia.org/wiki/OLED}
    \bibitem{bib:SSD}
    \url{https://en.wikipedia.org/wiki/Seven-segment_display}
    \bibitem{bib:Why} Ucl, S. A. . Shedding light on why blue LEDS are so tricky to make. Post on Jan 09, 2015 at \url{http://www.chemeurope.com/en/news/151134/}.
    \bibitem{bib:Why2} Chilton, Alexander. (2019, May 07). Nobel Prize in Physics 2014: Why Were Blue LEDs so Hard to Make?. AZoM. Retrieved on April 27, 2022 from \url{https://www.azom.com/article.aspx?ArticleID=11451}.
\end{thebibliography}
\end{frame}

\addtocounter{framenumber}{-1}
\begin{frame}[plain]
    \begin{center}
        \Huge\textbf{谢谢大家}
    \end{center}
\end{frame}

\end{document}
